\documentclass[11pt]{article}
\usepackage{mathblatt}
% Präambel des Setzers (werkzeuge/setzer-praeambel.tex, wortgleich); dazu \gteil (Muster T2U)
\makeatletter
\setlength{\mbpfbe}{0mm}% keine Punkte (pruefung.md § 5: nur beim Prüfstein)
\newlength{\gEin}\setlength{\gEin}{0pt}
\let\gSchrift\relax
\renewcommand{\pfkreuzzeile}[1]{\par\noindent\hangindent3.2em\hangafter1\hspace*{1.6em}$\square$\ #1\par}
\newcommand{\gkreuz}[1]{$\square$\ #1\hspace{2em}}
\newcommand{\gkreuzzeile}[2]{\par\noindent\hangindent2.2em\hangafter1\mbox{#1\,$\square$}\ #2\par}% Buchstabe, Option: eine je Zeile (neben Skizze)
% links, was man ansieht, rechts, was man tut; Spaltengrenze überall an derselben Stelle (gemeinsam.md „Satz“)
\newsavebox{\gbox}\newlength{\gLinks}
% v0.10 (Lehrer 13.10.): keine starre Halbierung – die linke Spalte ist so breit wie die Figur (+ 6 mm), die
% Antwort steht direkt daneben; ist die Figur breiter als 0,62 Zeile, steht die Antwort darunter
\newcommand{\gzwei}[2]{\par\smallskip\noindent\sbox{\gbox}{#1}%
  \setlength{\gLinks}{\dimexpr\wd\gbox+6mm\relax}%
  \ifdim\gLinks>0.62\linewidth
    \usebox{\gbox}\par\smallskip\noindent\begin{minipage}{\linewidth}#2\end{minipage}%
  \else
    \begin{minipage}[t]{\gLinks}\vspace{0pt}\usebox{\gbox}\end{minipage}%
    \begin{minipage}[t]{\dimexpr\linewidth-\gLinks\relax}\vspace{0pt}\raggedright #2\end{minipage}%
  \fi\par}
% v0.9 (12.10., pruefung.md § 5 „Marke mit B, Bedingungen rechts“): drittes Argument von pfaufg = Bedingungen
% der Prüfung (v0.10: \pfohnetr = TR durchgestrichen, „FOR“), klein grau im rechten Rand auf der Grundlinie der Nummer, Zeilen mit \\
\newcommand{\gbedrand}[2]{\ifblank{#2}{}{\rlap{\hspace*{\dimexpr#1+2mm\relax}%
  \parbox[t]{15mm}{\raggedright\scriptsize\leavevmode\color{mbgrau}\strut#2}}}}
% Teil eines Bündels (a, b …): Marke links im Rand, Bedingungen rechts im Rand, auf der Zeile des Teils
\newcommand{\gteilrand}[2]{\llap{\mbpfrandmarke{#1}\hspace*{\mbpfnr}}\gbedrand{\linewidth}{#2}}
% Teil eines Bündels: Buchstabe in eigener Spalte, Inhalt eingerückt (\gEin = Einzug, damit die Spaltengrenze von
% \gzwei an derselben Stelle bleibt)
\newenvironment{gbteil}[3]{\par\smallskip\noindent\gteilrand{#2}{#3}\makebox[\mbpfteil][l]{\textbf{#1}}%
  \begin{minipage}[t]{\dimexpr\linewidth-\mbpfteil\relax}\setlength{\gEin}{\mbpfteil}\raggedright\strut\ignorespaces}%
 {\end{minipage}\par}
% Grafik höchstens so breit wie die Spalte (Bündel nebeneinander, Koordinatensystem ohne Karo)
\newsavebox{\gpbox}
\newcommand{\gpassend}[1]{\sbox{\gpbox}{#1}\ifdim\wd\gpbox>\linewidth\typeout{GPASSEND: verkleinert \the\wd\gpbox\space auf \the\linewidth}\resizebox{\linewidth}{!}{\usebox{\gpbox}}\else\usebox{\gpbox}\fi}
% Skizzenfeld (v0.10, Lehrer 13.10.): Karo ab dem linken Rand über die Zeile, „Skizze“ klein grau oben links im Feld
% (wie \pfskizzenfeld der Vorlage); der rechte Rand bleibt den Bedingungen
\newcommand{\gskizze}[1]{\par\smallskip\noindent\begin{tikzpicture}
  \clip (0,0) rectangle ({\linewidth-1mm},{#1*10mm});
  \draw[black!16,line width=0.3pt,step=5mm] (0,0) grid ({\linewidth},{#1*10mm});
  \node[anchor=north west,inner sep=2pt,font=\scriptsize,text=mbgrau,fill=white] at (0,{#1*10mm}) {Skizze};
  \end{tikzpicture}\par}
\RenewDocumentEnvironment{pfaufg}{m m m}{%
  \begin{lrbox}{\mbaufgabebox}\begin{minipage}[t]{\linewidth}%
  \gSchrift\noindent\hspace*{\gEin}\mbpfrandmarke{#1}\gbedrand{\linewidth-\gEin}{#3}%
  \begin{minipage}[t]{\mbpfnr}\strut\textbf{#2}\end{minipage}%
  \begin{minipage}[t]{\dimexpr\linewidth-\gEin-\mbpfnr\relax}\raggedright\strut\ignorespaces}%
 {\end{minipage}%
  \end{minipage}\end{lrbox}%
  \setlength{\mbaufgabehoehe}{\dimexpr\ht\mbaufgabebox+\dp\mbaufgabebox\relax}%
  \ifdefempty{\gkopfnach}{\par\addvspace{5pt}\Needspace*{\mbaufgabehoehe}}{%
    % Stufenkopf, der mit dieser (ersten) Aufgabe zusammenbleibt: Platz für beide verlangen, dann den Kopf setzen
    % (12.10., Quadratische: ein Kopf stand allein unten, die Aufgabe mit Koordinatensystem auf der nächsten Seite)
    \par\Needspace*{\dimexpr\mbaufgabehoehe+40pt\relax}\gkopfnach\gdef\gkopfnach{}\par\addvspace{5pt}}%
  \noindent\usebox{\mbaufgabebox}\par\addvspace{5pt}}
% Stufenkopf vormerken: \gkopfvor{<Kopf>} – die nächste pfaufg setzt ihn, sobald beide zusammen Platz haben
\def\gkopfnach{}
\newcommand{\gkopfvor}[1]{\gdef\gkopfnach{#1}}
% Teilaufgabe mit Skizze links in gzwei: Buchstabe in eigener Spalte, darunter Text und Skizze
\newcommand{\gteil}[2]{\makebox[\mbpfteil][l]{\textbf{#1}}\parbox[t]{\dimexpr0.5\textwidth-\gEin-\mbpfnr-5mm-\mbpfteil\relax}{\raggedright #2}}
% Stufenkopf der Vorlage (mathblatt.sty 2026-10-12: Linie, Abstand, 13 pt, \pfjahre) ohne ihren Vorhalt von
% 14 Zeilen: der Setzer fordert vor dem Kopf 0,16 Seitenhöhe (Kopf + erste Aufgabe), mehr nicht
\setlength{\mbpfstufevorhalt}{0pt}
\makeatother
\newcommand{\karo}[1]{\par\smallskip\noindent\begin{tikzpicture}
  \clip (0,0) rectangle ({\linewidth-1mm},{#1*10mm});
  \draw[black!16,line width=0.3pt,step=5mm] (0,0) grid ({\linewidth},{#1*10mm});
  \end{tikzpicture}\par}
\newcommand{\kk}[1]{$\square$\,#1\hspace{0.9em}}
\newcommand{\linie}[1]{\,{\color{black!45}\rule[-1pt]{#1}{0.4pt}}\,}
\newcommand{\pkt}{\textperiodcentered{}}
% v0.14 (14.10.): Stufenkopf mit fünf Punkten (Option = Jahre für die Bestellkarte) oder „zuletzt 19“
\NewDocumentCommand{\gpunkte}{o m}{#2}
\newcommand{\gzuletzt}[1]{zuletzt #1}
% v0.7 (12.10.): Schülerfrage unter dem Stufenkopf; Ankreuzoption als Bild (Buchstabe und Kästchen darunter)
\newcommand{\gfrage}[1]{\par\nobreak\noindent\hspace*{\mbpfnr}{\small\itshape #1}\par\nobreak\addvspace{2pt}}
\newlength{\gbildbreite}\setlength{\gbildbreite}{0.32\linewidth}% Setzer setzt sie je Reihe (vier Bilder: 0,24)
\newcommand{\gbildoption}[3]{\begin{minipage}[t][#3][b]{\gbildbreite}\centering #2\par\smallskip\mbox{#1\,$\square$}\end{minipage}}
% Lineare (12.10.): Bild mit Buchstabe ohne Kästchen (Diagramm zuordnen); Aussage mit „wahr / falsch“ zum Ankreuzen
\newcommand{\gbildname}[3]{\begin{minipage}[t][#3][b]{\gbildbreite}\centering #2\par\smallskip\mbox{\textbf{#1}}\end{minipage}}
\newcommand{\gwahrfalsch}[2]{\par\noindent\hangindent2.2em\hangafter1\makebox[2.2em][l]{(#1)}#2\hfill\mbox{wahr $\square$\quad falsch $\square$}\par}
\begin{document}

\pfheftstil
\pflaufkopf[Abitur GK \pkt{} Prüfung]{Nullstellen}{AA20}\pfseitevon
\pfheftkopf{Nullstellen}{}
\pfstufenzeile{Funktionswert gesucht \pkt{} Nullstellen gesucht \pkt{} Schnittpunkte mit den Achsen gesucht}

% ===================== Stufe: Funktionswert gesucht =====================
\gkopfvor{\pfstufekopf{Funktionswert gesucht}{\gpunkte[26 25]{\textcolor{black!80}{$\bullet$}\,\textcolor{black!80}{$\bullet$}\,\textcolor{black!25}{$\circ$}\,\textcolor{black!25}{$\circ$}\,\textcolor{black!25}{$\circ$}}}\special{pdf:dest (stufe-funktionswert-gesucht) [@thispage /XYZ @xpos @ypos null]}}
% kein Einstieg: keine Schülerfrage in der Gliederung (Schülerfrage, -Fälle, -Form, -Weg, -Skizze; oder Frage) – Einstieg nur von dort, nie aus der Bank (pruefung.md § 5, 12.10.)
% 2026-bb-gk-B2.2g-m1 (Original 2026-bb-gk-B2.2g, Schwere I, rang 1; herausgelöste Fassung: keine)
\begin{pfaufg}{GK~’26\\B2.2g}{1.}{}
Ein Speicherofen kühlt ab. Die Temperatur seiner Oberfläche in °C wird beschrieben durch $k(x) = 55e^{-\tfrac{x}{300}} + 18$; $x$ ist die Zeit in Minuten seit Beginn des Abkühlens. Gib die Temperatur zu Beginn an.
\hfill\linie{30mm}
\fusshilfe{\mbox{1:~73 °C}}
\end{pfaufg}
% 2025-bebb-gk-A1.4a-m1 (Original 2025-bebb-gk-A1.4a, Schwere I, rang 2; herausgelöste Fassung: keine)
\begin{pfaufg}{GK~’25\\A1.4a}{2.}{\pfohnetr}
Gegeben ist die in $\mathbb{R}$ definierte Funktion $f$ mit $f(x) = 2e^x - 2e^{-1}$; die Abbildung zeigt ihren Graphen. Weise nach, dass $-1$ eine Nullstelle von $f$ ist.
\gzwei{\begin{abiksys}[xmin=-4,xmax=1,ymin=-1,ymax=5,le=0.5]\funktion[0.6]{2*exp(\x)-2*exp(-1)}{G_f}\end{abiksys}}{\karo{1}}
\fusshilfe{\mbox{2:~0}}
\end{pfaufg}

% ===================== Stufe: Nullstellen gesucht =====================
\gkopfvor{\pfstufekopf{Nullstellen gesucht}{\gpunkte[24 23 22]{\textcolor{black!25}{$\circ$}\,\textcolor{black!25}{$\circ$}\,\textcolor{black!80}{$\bullet$}\,\textcolor{black!80}{$\bullet$}\,\textcolor{black!80}{$\bullet$}}}\special{pdf:dest (stufe-nullstellen-gesucht) [@thispage /XYZ @xpos @ypos null]}}
% kein Einstieg: keine Schülerfrage in der Gliederung (Schülerfrage, -Fälle, -Form, -Weg, -Skizze; oder Frage) – Einstieg nur von dort, nie aus der Bank (pruefung.md § 5, 12.10.)
% 2022-bebb-gk-A1.1a-m1 (Original 2022-bebb-gk-A1.1a, Schwere I, rang 3; herausgelöste Fassung: keine)
\begin{pfaufg}{GK~’22\\A1.1a}{3.}{\pfohnetr}
Berechne die Nullstellen der Funktion $f$ mit $f(x) = x^3 - 10x^2 + 21x$.
\karo{2}
\fusshilfe{\mbox{3:~Nullstellen 0, 3 und 7}}
\end{pfaufg}
% 2023-bebb-gk-B2.2a-m1 (Original 2023-bebb-gk-B2.2a, Schwere II, rang 4; herausgelöste Fassung: keine)
\begin{pfaufg}{GK~’23\\B2.2a}{4.}{}
Gegeben ist die Funktion $f$ mit $f(x) = x^3 - 9x^2 + 24x - 16$, $x \in \mathbb{R}$. Der Term lässt sich auch schreiben als $f(x) = (x - 4) \cdot (x^2 - 5x + 4)$. Gib die Koordinaten des Punkts an, in dem der Graph von $f$ die $y$-Achse schneidet. Ermittle, wie viele Nullstellen $f$ hat.
\karo{2}
\fusshilfe{\mbox{4:~Punkt (0|$-$16); genau zwei Nullstellen, 1 und 4}}
\end{pfaufg}
% 2024-bebb-gk-B2.2b-m1 (Original 2024-bebb-gk-B2.2b, Schwere II, rang 5; herausgelöste Fassung: keine)
\begin{pfaufg}{GK~’24\\B2.2b}{5.}{}
Gegeben ist die Funktion $f$ mit $f(x) = 0{,}5x^4 - 3x^2 + 2{,}5$, $x \in \mathbb{R}$; ihr Graph ist $G_f$. Der Punkt $(1 \mid 0)$ liegt auf $G_f$. Begründe ohne Rechnung, dass auch $(-1 \mid 0)$ auf $G_f$ liegt. Bestimme die übrigen Schnittpunkte von $G_f$ mit der $x$-Achse.
\karo{2}
\fusshilfe{\mbox{5:~($-$1|0) liegt auf $G_f$; ($-$√5|0) und (√5|0)}}
\end{pfaufg}

% ===================== Stufe: Schnittpunkte mit den Achsen gesucht =====================
\gkopfvor{\pfstufekopf{Schnittpunkte mit den Achsen gesucht}{\gpunkte[26 24 23 22]{\textcolor{black!80}{$\bullet$}\,\textcolor{black!25}{$\circ$}\,\textcolor{black!80}{$\bullet$}\,\textcolor{black!80}{$\bullet$}\,\textcolor{black!80}{$\bullet$}}}\special{pdf:dest (stufe-schnittpunkte-mit-den-achsen-gesucht) [@thispage /XYZ @xpos @ypos null]}}
% kein Einstieg: keine Schülerfrage in der Gliederung (Schülerfrage, -Fälle, -Form, -Weg, -Skizze; oder Frage) – Einstieg nur von dort, nie aus der Bank (pruefung.md § 5, 12.10.)
% 2026-bb-gk-B2.2a-m1 (Original 2026-bb-gk-B2.2a, Schwere I, rang 6; herausgelöste Fassung: keine)
\begin{pfaufg}{GK~’26\\B2.2a}{6.}{}
Gegeben ist die Funktion $f$ mit $f(x) = 2e^x + 3$, $x \in \mathbb{R}$; ihr Graph ist $G_f$. Begründe am Funktionsterm, dass $G_f$ die $x$-Achse nicht schneidet. Gib die Wertemenge von $f$ an.
\karo{2}
\fusshilfe{\mbox{6:~keine Nullstelle; alle y mit y > 3}}
\end{pfaufg}
% 2024-bebb-gk-B2.1a-m1 (Original 2024-bebb-gk-B2.1a, Schwere I, rang 7; herausgelöste Fassung: keine)
\begin{pfaufg}{GK~’24\\B2.1a}{7.}{}
Gegeben ist die in $\mathbb{R}$ definierte Funktion $f$ mit $f(x) = (x - 3) \cdot e^{-\tfrac{x}{3} + 2}$ und ihrer Ableitung $f'(x) = \left(2 - \tfrac{x}{3}\right) \cdot e^{-\tfrac{x}{3} + 2}$. Die Abbildung zeigt $G_f$ und $G_{f'}$ ohne Skalierung. $f$ hat genau eine Nullstelle. Gib sie an. Ermittle die Koordinaten des Punkts, in dem $G_f$ die $y$-Achse schneidet.
\gzwei{\begin{abiksys}[xmin=-1,xmax=10,ymin=-24,ymax=16,le=0.5,yle=0.06,ystep=5,ohnezahlen,ohnegitter]\funktion[1.5]{(\x-3)*exp(-\x/3+2)}{G_f}\funktion[1]{(2-\x/3)*exp(-\x/3+2)}{G_{f'}}\end{abiksys}}{\karo{2}}
\fusshilfe{\mbox{7:~3; Punkt (0|$-$3e²)}}
\end{pfaufg}
% 2022-bebb-gk-B2.1a-m1 (Original 2022-bebb-gk-B2.1a, Schwere I, rang 8; herausgelöste Fassung: keine)
% Zweitthema ausgeblendet (Spalten <feld>_einheit, pruefheft-p.md v0.9)
\begin{pfaufg}{GK~’22\\B2.1a}{8.}{}
Gegeben ist die in $\mathbb{R}$ definierte Funktion $f$ mit $f(x) = (x + 3) \cdot e^{-x}$. Gib die Punkte an, in denen der Graph von $f$ die Koordinatenachsen schneidet.
\karo{2}
\fusshilfe{\mbox{8:~Schnittpunkte ($-$3|0) und (0|3)}}
\end{pfaufg}
% 2023-bebb-gk-B2.1b-m1 (Original 2023-bebb-gk-B2.1b, Schwere II, rang 9; herausgelöste Fassung: keine)
\begin{pfaufg}{GK~’23\\B2.1b}{9.}{}
Gegeben ist die in $\mathbb{R}$ definierte Funktion $f$ mit $f(x) = 0{,}5 \cdot (x^2 - 9) \cdot e^x$; ihr Graph ist $G$. Bestimme die Koordinaten der Punkte, in denen $G$ die beiden Koordinatenachsen schneidet. Diese drei Punkte bilden zusammen mit einem Punkt $R$ die Ecken einer Raute. Gib die Koordinaten von $R$ an.
\karo{2}
\fusshilfe{\mbox{9:~Schnittpunkte ($-$3|0), (3|0) und (0|$-$4,5); R(0|4,5)}}
\end{pfaufg}

\pfblattende{10}{26 & 25 & 24 & 23 & 22 & 21 & 20 & 19 & 18 & 17}{B2.2a B2.2g & A1.4a & B2.1a B2.2b & B2.1b B2.2a & A1.1a B2.1a & B2.1a B2.1b B2.2a & B2.1a B2.2a & B2.1a &  & B1.2a}%
{\pfausschnitt{Grenzverhalten und Symmetrie}{Nullstellen}{Extrempunkte \pkt{} Monotonie}}
\end{document}